\documentclass[11pt]{article}
\usepackage[a4paper,margin=25mm]{geometry}
\usepackage[T1]{fontenc}
\usepackage[utf8]{inputenc}
\usepackage{amsmath,amssymb,booktabs}
\usepackage{enumitem}
\usepackage{array}
\usepackage{microtype}
\usepackage{xurl}
\usepackage[hidelinks]{hyperref}

\setlength{\parindent}{0pt}
\setlength{\parskip}{0.55em}
\setlist{nosep,leftmargin=*}
\emergencystretch=3em
\newcolumntype{L}[1]{>{\raggedright\arraybackslash}p{#1}}

\title{M13-C3 --- H4 cautious vs brave\\
roots \(a\)/\(b\) on the closed H0 table}
\author{}
\date{H4 run / analysed / not a Bucket~3 claim}

\begin{document}
\maketitle

\section{Status}

\textbf{H4 run / analysed.}
Repository \texttt{baseline\_cautious} vs locked ECAI brave on fixture
\path{m13_bucket3_binary_collider_and} / \texttt{n30\_seed42} is complete.
Markdown twin: \texttt{h4\_cautious\_vs\_brave.md}.

\textbf{Lead finding:} switching to cautious \textbf{blocks the brave
residual-assumption gadget} that had allowed roots \(a\)/\(b\) to
\texttt{solved}; paths coincide through early \(\alpha_1\)--\(\alpha_3\)
setup, then brave says \textbf{OK} and finishes while cautious says
\textbf{KO} and returns \texttt{completed\_no\_solution}.
\textbf{Control:} child \(c\) still solves with a byte-identical delta to
locked ECAI \(c\).
H0--H3 remain closed.
\textbf{Not a Bucket~3 claim.}
H4b analysed separately (\texttt{h4b\_cautious\_split\_under\_a.md});
H5 and H6 analysed separately (\texttt{h5\_greedy\_cautious.md},
\texttt{h6\_folding\_and\_n\_ablation.md}); H7 signposted.

\section{Research question}

On the closed H0 AND-collider table, does repository-baseline cautious
learning change the ECAI-brave root outcomes for \(a\)/\(b\) that relied on a
residual per-row assumption gadget (rows~9 vs~26), and what happens to
control \(c\)?

\section{Main finding (lead)}

Locked ECAI brave solves roots \(a\)/\(b\) with assumption-rich
\(\alpha\)-nests.
Repository \texttt{baseline\_cautious} on the \textbf{same} table returns
\texttt{completed\_no\_solution} for both roots (no \texttt{delta.aba}).

\begin{center}
\begin{tabular}{@{}lL{0.38\textwidth}L{0.42\textwidth}@{}}
\toprule
Target & Locked ECAI brave & \texttt{baseline\_cautious} \\
\midrule
\(a\) & \texttt{solved} (\(\alpha\)-nest) &
  \texttt{completed\_no\_solution} ($\sim$70s; $\sim$5k-line stdout) \\
\(b\) & \texttt{solved} (mirror) &
  \texttt{completed\_no\_solution} ($\sim$66s) \\
\(c\) (control) & \texttt{solved} &
  \texttt{solved} (byte-identical delta; $\sim$1.3s; audit SAT) \\
\bottomrule
\end{tabular}
\end{center}

\textbf{Learning-mode / entailment semantics change the root outcome} by
rejecting the 9-vs-26-style per-row assumption cover.
This is \textbf{not} ``cautious never solves'' (see control \(c\)).

\section{Decisive procedural fork}

Both arms share the early path through \(\alpha_1\)--\(\alpha_3\).
Table conflict: row~9 is \(E^+\) \((1,0,0)\); row~26 is \(E^-\) \((0,0,0)\) ---
identical BK \texttt{b\_val\_0},\texttt{c\_val\_0}, opposite labels.

\textbf{Brave:}
\begin{verbatim}
folding result: c_alpha_3(A) <- [b_val_0(A)]
 found: alpha_2/1 ... OK, assumption introduction result:
   c_alpha_3(A) <- [alpha_2(A),b_val_0(A)]
\end{verbatim}
Closing nest:
\begin{verbatim}
a(A) :- alpha_2(A), b_val_0(A).
c_alpha_2(A) :- alpha_3(A), c_val_0(A).
c_alpha_3(A) :- alpha_2(A), b_val_0(A).
\end{verbatim}

\textbf{Cautious:}
\begin{verbatim}
folding result: c_alpha_3(A) <- [b_val_0(A)]
 found: alpha_2/1 ... KO, cannot introduce an assumption for [b_val_0(A)]
\end{verbatim}
then budget exhaustion and \texttt{* No solution found!}.
Target \(b\) mirrors.
Paths coincide until the brave gadget; semantics then diverge.

\section{Control \(c\)}

Cautious \(c\) still solves with a byte-identical delta to locked ECAI \(c\):
\begin{verbatim}
c(A) :- alpha_1(A), a_val_1(A).
c_alpha_1(A) :- b_val_0(A).
assumption(alpha_1(A)).
contrary(alpha_1(A),c_alpha_1(A)) :- assumption(alpha_1(A)).
\end{verbatim}
Mode change is selective.
String match / \texttt{solved} is not claimed as causal recovery.

\section{Setup / paths}

Fixture/sample: closed H0 \path{m13_bucket3_binary_collider_and} /
\texttt{n30\_seed42}
(sha256 \texttt{1edae465\ldots a8ca}).
Brave arm: locked \texttt{ecai2024}.
Cautious arm: \texttt{baseline\_cautious}
(\texttt{configs/baseline\_cautious\_config.pl}; not ECAI/AAMAS/H5).
Primary targets \(a\),\,\(b\); control \(c\).

\section{Semantic boundary}

Authoritative for acceptance: Prolog learner outcome under the selected
\texttt{learning\_mode}.
Under cautious learning, SAT of \texttt{bk.sol\_chk.asp} is an existential
final-artefact integrity witness, not proof that every positive is a cautious
consequence.

\section{Boundaries / non-claims}

\begin{itemize}
  \item Not a Bucket~3 claim.
  \item Cautious root no-solution is not mechanism recovery.
  \item Brave root \texttt{solved} is not mechanism recovery.
  \item Control \(c\) string match is not causal recovery.
  \item Not a ranking of brave vs cautious as generally better.
  \item H5 and H6 analysed separately (\texttt{h5\_greedy\_cautious.md},
    \texttt{h6\_folding\_and\_n\_ablation.md}); H7 signposted only.
    H6/H7 are not this H4 evidence body.
\end{itemize}

\section{Prompted follow-ups (signpost only)}

H5 is now run / analysed separately (\texttt{h5\_greedy\_cautious.md}).
H6 (length drivers: \(n\), \texttt{folding\_steps} under fixed cautious) and
H7 (brave \texttt{asm\_intro(sechk)} vs \texttt{relto} on the 9/26 conflict)
remain future probes prompted by H4 traces, not yet approved runs.

\section{Next}

H4 documentation is complete.
H4b and H5 are run / analysed separately.
H6 is analysed separately; H7 is signposted.
Still \textbf{no Bucket~3 claim}.

\end{document}
